\section{Обсуждение и выводы}
\label{sec:discussion}

\subsection{Рабочая рекомендация и её цена}
\label{sec:disc-recommendation}

Расчётный оптимум лежит при $g\in[0.5,\,0.75]\ \text{нм}$
(раздел~\ref{sec:res-threshold}). Рекомендовать именно его, однако, было бы
неосторожно: при субнанометровых зазорах существенны три эффекта, которых в
принятой модели нет. Перенос заряда в металлических зазорах включается ниже
$\approx0.3$--$\SI{0.5}{\nano\meter}$; нелокальное экранирование становится
заметным на том же масштабе; наконец, лигандная оболочка коллоидной КТ обычно
составляет $0.5$--$\SI{2}{\nano\meter}$, так что зазор $\SI{0.5}{\nano\meter}$
требует специальной замены лигандов и в стандартной коллоидной химии
недостижим.

Цена отступления в защитимую область известна точно, и она невелика именно
благодаря насыщению~\eqref{eq:saturation}:
\begin{equation}
  \frac{\mathcal F_\eta(\SI{1.0}{\nano\meter})}{\mathcal F_\eta(\SI{0.5}{\nano\meter})} = 1.20,
  \qquad
  \frac{\mathcal F_\eta(\SI{1.5}{\nano\meter})}{\mathcal F_\eta(\SI{0.5}{\nano\meter})} = 1.44 .
  \label{eq:step-back-cost}
\end{equation}
Поэтому мы сообщаем расчётный оптимум как результат модели, а в качестве
рабочей рекомендации предлагаем интервал
\begin{equation}
  g \in [1.0,\;1.5]\ \text{нм},
  \label{eq:recommendation}
\end{equation}
теряя $20$--$44\,\%$ по требуемому флюенсу и приобретая область, где все
допущения модели защитимы, а конфигурация собирается обычными средствами.

Этот интервал удачен и по второй наблюдаемой. Слабополевое усиление на
фиксированной частоте имеет максимум $G_{\rm exc}=25.5$ при
$\SI{1.5}{\nano\meter}$, а при $\SI{1.0}{\nano\meter}$ равно $25.3$ — в пределах
$1\,\%$ от максимума. То есть рекомендованное окно практически оптимально для
узкополосной схемы и стоит умеренной потери для импульсной.

\subsection{Что выбирать в первую очередь}
\label{sec:disc-geometry}

Главный практический вывод работы касается не расстояния. Выбор канала — то
есть места КТ на частице и ориентации поля — меняет порог от выигрыша в $45.7$
раза до конфигураций, где порог вообще не достигается, тогда как весь диапазон
зазоров внутри лучшего канала даёт множитель $13.3$. Правило при этом простое и
проверяется по табл.~\ref{tab:thresholds}: усиление возникает, когда внешнее
поле \emph{нормально} к поверхности частицы в точке расположения КТ, и
сменяется подавлением, когда поле касательно. Иными словами, геометрию и
поляризацию следует фиксировать прежде, чем оптимизировать зазор.

Второй практический вывод: рекомендация зависит от схемы возбуждения, и эта
зависимость количественная, а не качественная. При узкополосном возбуждении на
фиксированной частоте оптимум внутренний, вблизи $\SI{1.5}{\nano\meter}$: сближение
одновременно усиливает поле и уводит гибридный резонанс с частоты зонда, и при
$\SI{0.5}{\nano\meter}$ смещение достигает половины ширины линии. Если же лазер
перестраивается вслед за резонансом либо возбуждение широкополосное
(\SI{20}{\femto\second}, спектральная ширина \SI{91}{\milli\electronvolt} против
смещения $\lesssim\SI{1.4}{\milli\electronvolt}$), расстройка не стоит ничего и
оптимума нет — отклик монотонно улучшается до контакта. Перед оптимизацией
геометрии нужно, следовательно, зафиксировать схему измерения.

\subsection{Когда допустимы упрощённые модели}
\label{sec:disc-models}

Методический результат работы отрицательный для обоих упрощений, и в обоих
случаях он количественный.

\paragraph{Дипольное приближение.} Расхождение с точным решением убывает с
расстоянием как $R^{-3}$, но объявленного допуска $10\,\%$ в допустимой
квазистатической области не достигает ни для одного из двух каналов с
наибольшим усилением. Причина установлена экстраполяцией измеренного закона
и выражена неравенством~\eqref{eq:no-window}: нужная точность достигалась бы при
$R\approx\SI{34}{\nano\meter}$, тогда как квазистатика держится до
$R\le\SI{29.4}{\nano\meter}$. Иными словами, точечно-дипольная замена частицы
становится достаточно точной только там, где обе модели уже неприменимы по
другой причине.

Это не означает, что дипольным приближением нельзя пользоваться никогда: для
трёх слабых каналов граница нашлась и составляет $7$--$\SI{10}{\nano\meter}$.
Вывод формулируется избирательно: \emph{дешёвая модель отказывает именно в том
режиме, ради которого систему и собирают}. Причём она ошибается не только
количественно (завышение требуемого флюенса в $2.4$ раза при
$\SI{1}{\nano\meter}$), но и качественно: для канала «вершина, поперёк поля» она
предсказывает конечный порог при всех зазорах, тогда как точное решение
показывает, что при $g\le\SI{3}{\nano\meter}$ порог недостижим.

\paragraph{Одноосцилляторная дисперсия.} Ошибка линейного описания материала
($\mathrm{NRMS}=0.72$ против $0.025$) переносится в нелинейный режим как
завышение порога в $7.7$ раза, а расхождение населённостей достигает $0.87$ —
почти всего доступного диапазона. Для золота в окне
$0.8$--$\SI{3.5}{\electronvolt}$, где сосуществуют друдевский отклик и межзонные
переходы, одного лоренцева контура недостаточно. Это существенно для работ,
переносящих одноосцилляторную модель из задач о благородных металлах в узкой
полосе на короткоимпульсные задачи.

\subsection{Границы полученных выводов}
\label{sec:disc-limits}

Абсолютные значения усилений и порогов несут систематическую неопределённость
эффектов, оценённых, но не включённых в модель
(раздел~\ref{sec:applicability}): запаздывание даёт множитель $1.29$ для
$|\alpha|^2$ на несущей и смещение плазмонного резонанса на
$-\SI{21}{\milli\electronvolt}$, поверхностное затухание — множитель
$0.83$--$0.95$. Обе оси сравнения работы при этом свободны от этих поправок:
и DD против FQS, и $N=1$ против $N=12$ используют одну и ту же геометрию и один
и тот же материал, так что поправки сокращаются в отношениях.

Отдельно следует подчеркнуть, что точечное приближение КТ нарушено во всём
рассмотренном диапазоне, а не только при малых зазорах: индикатор
$r_{\rm QD}|\nabla E|/|E|\ge3r_{\rm QD}/R$ меняется лишь от $0.35$ до $0.32$
между $0.5$ и $\SI{1.5}{\nano\meter}$, поскольку расстояние до центра частицы
определяется её большой полуосью. Отступление в рекомендованный
интервал~\eqref{eq:recommendation} устраняет туннелирование и нелокальность, но
не конечность размера КТ. Усреднение поля по объёму точки — естественное
направление уточнения; оно изменит абсолютные величины, но не отношения, на
которых построены оба методических вывода.

Наконец, результаты относятся к одной частице заданных размеров, одному
материалу и одной длительности импульса; ранжирование выполнено по
сконфигурированному дискретному набору конфигураций и не является глобальным
оптимумом по форме и размеру частицы.

\subsection{Выводы}
\label{sec:disc-conclusions}

\begin{enumerate}
  \item Наибольшее снижение порога импульсного возбуждения даёт КТ у вершины
    вытянутого сфероида при поле вдоль оси: в $45.7$ раза относительно
    изолированной КТ. Расчётный оптимум по зазору — интервал
    $[0.5,\,0.75]\ \text{нм}$; в качестве рабочей рекомендации предлагается
    $[1.0,\,1.5]\ \text{нм}$ ценой $20$--$44\,\%$ по флюенсу.
  \item Определяющим является не зазор, а ориентация поля относительно
    поверхности: нормальная даёт усиление, касательная — подавление вплоть до
    недостижимости порога.
  \item Оптимальный зазор зависит от схемы возбуждения: для узкополосного
    зондирования на фиксированной частоте он внутренний ($\approx\SI{1.5}{\nano\meter}$),
    для широкополосного импульса оптимума нет и отклик насыщается.
  \item Точечно-дипольная модель частицы не достигает точности $10\,\%$ для
    сильных каналов нигде внутри области применимости квазистатики и ошибается
    качественно; для слабых каналов её граница составляет
    $7$--$\SI{10}{\nano\meter}$.
  \item Одноосцилляторное описание дисперсии золота завышает порог в $7.7$ раза
    и для короткоимпульсных задач непригодно.
\end{enumerate}
